\section{The absolutely convergent kernel on the selected family}
Fix a datum $D$ and a spectral tie $T$ on it. This records a setting
$X$, full arithmetic data $Y$, and, for every proof $h_F$ of the
finiteness record of $Y$, a comparison from the original spectral
datum $D.s$ to the constructed datum $s_Y=\mathrm{sOf}(h_F)$.
The expansion below is an identity on the original column spaces
of $D.s$, including its normalizing scalar.

% BEGIN KXE1
\paragraph{KXE1. The dimension in the normalizing constant.}
For any supplied setting $X$ and word $u$, the dimension of its
compact representation $V_u$ is the product $\mathrm{dimPi}(u)$
in the definition of the abstract Gram kernel.
\begin{proof}
At a place $v$ let $\lambda$ be the compact Harish--Chandra parameter
attached by $X.\mathrm{pl}$ to $u$. Its entries are strictly decreasing.
The compact Weyl dimension formula gives, in these shifted coordinates,
\[
            \dim V_{u,v}=
        \frac{\prod_{a<b}(\lambda_a-\lambda_b)}
             {\prod_{a=0}^{k-1}a!}.
\]
Indeed the positive roots are the differences of coordinates,
and their pairings with the half-sum of positive roots give
$\prod_{a<b}(b-a)=\prod_{a=0}^{k-1}a!$. The source defines its
local $\mathrm{dimPiU}$ as the least dimension of a representation
of this type. The given $V_{u,v}$ belongs to the defining set and
the same formula gives the same dimension for every member.
This proves the local identity.

Choose orthonormal bases $b_{v,j}$ and take the supplied tensor
vectors $e_{\mathbf j}=\mathrm{tens}_u((b_{v,j_v})_v)$.
The product inner-product axiom makes them orthonormal. For
$y=\mathrm{tens}_u((x_v)_v)$, local Parseval and the same axiom give
\[
 \sum_{\mathbf j}|\langle e_{\mathbf j},y\rangle|^2
 =\prod_v\sum_j|\langle b_{v,j},x_v\rangle|^2
 =\prod_v\|x_v\|^2=\|y\|^2 .
\]
Consequently $y$ belongs to their span, by orthogonal projection.
All the supplied tensors span $V_u$, so these particular ones form
a basis. Counting the multi-indices and using the given bijection
from places to the indices of $\mathrm{dimPi}$ gives the assertion.
This proof uses the inner-product and spanning fields even when the
tensor map is presented without multilinearity axioms.
\end{proof}
% END KXE1

% BEGIN KXE2
\paragraph{KXE2. The original kernel equals the sum of its terms.}
Assume the supplied $X=T.X$ satisfies $\mathrm{WuDecay}$ and its
archimedean kernel is integrable. For all ordered word pairs $(u,u')$,
\[
 K_{D.s}(u,u')=\sum_t A_t,\qquad
                  \sum_t\|A_t\|<\infty ,
 \qquad A_t=(D.s).\mathrm{aTerm}(u,u',t).
\]
The operator sum converges as a net over finite subsets. This is
$\mathrm{PX.Exp}(D)$. The source gap on $D$ suffices for both
analytic assumptions.
\begin{proof}
SFS2 gives the level normalization for these $X,Y$. In the proof of
QPR11 the three hypotheses extracted from the concrete tie were
that normalization, the decay of the supplied wave and integrability
of $X.\mathrm{ph}$. To see their sufficiency on the given data,
FLD3 folds both word columns, and QPR3 places their period series
in the same Hilbert range. QPR4 chooses strict domains, QPR5--QPR8
give the orbital integrals and the absolutely summable norm bounds,
and QPR9--QPR10 give independence of the domains and column
representatives. PGR18 then gives the ungrouped operator series with
its normalized Gram pairing for $Y$. This argument works for equal
words as well. SFS5 now supplies $h_F:\mathrm{FinX}(Y)$.
Use KXE1 at each of the two words and apply PGR20. Its finite-fibre
regrouping gives the asserted sum and absolute convergence first
on $s_Y$.

Take $I=T.hI(h_F)$, writing $U_w=I.eCu(w)$ and $U_H=I.eHH$
for its unitary identifications. Conjugation
$A\mapsto U_u A U_{u'}^{-1}$ preserves operator norm and finite
sums. The term comparison in $I$ sends each $A_t$ to the term on
$s_Y$ indexed by $I.eT(t)$, a bijection for this pair of words.
The period-map comparison is $U_HP_w=P'_w U_w$. Taking adjoints
gives
\[
 U_u P_u^*P_{u'}U_{u'}^{-1}=(P'_u)^*P'_{u'}.
\]
For $s=D.s$ or $s=s_Y$, write
\[
 c_w=s.\mathrm{volN}(w)s.\mathrm{volH}(w)/\mathrm{dimPi}(w).
\]
The volumes agree across $I$ and the parameters in $\mathrm{dimPi}$
are shared. Thus $(c_uc_{u'})^{-1/2}$, the real scalar multiplying
the operator product, agrees as well. The original $K_{D.s}$ maps
to $K_{s_Y}$. Pull the convergent finite sums back by inverse
conjugation and reindex their finite subsets by the inverse bijection;
the norm series pulls back with the same terms. This proves both
assertions on $D.s$. At a datum with the original gap SFS4 gives
both analytic hypotheses on $T.X$ itself.
\end{proof}
% END KXE2

% BEGIN KXE3
\paragraph{KXE3. One threshold for the kernel expansion at every index.}
An original family with a supplied $\mathrm{SSelF}$ has a single
threshold above which $\mathrm{PX.Exp}(D_{x,i})$ holds for
\emph{every} index $i$. In particular a class with a supplied
$\mathrm{SSel}$ satisfies the original closed $\mathrm{PX.ExpF}$.
\begin{proof}
Let $a$ be the threshold of the supplied selector, $b$ the threshold
for the gap in TFR2, and $x_1$ the strict family cutoff. At each
$x\ge\max(a,b,x_1+1)$ and each $i$, use the original $T_{x,i}$
and its own comparison. TFR2 supplies the gap at this index;
KXE2 supplies $\mathrm{PX.Exp}(D_{x,i})$ for every pair of words.
This choice of threshold works simultaneously for all the parameters
$C_0,C_g,c_b,C_u,c_T,C_T$ in $\mathrm{PX.Under}$. In particular
it works when its standing hypotheses hold with the family's own
$C',C_l$, giving precisely $\mathrm{ExpF}$. The argument on a class
takes the selector at the given family and its class-membership proof.
It works for every family of the class, so any of the outer closed
size thresholds for the identically true predicate suffice here.
\end{proof}
% END KXE3
